\documentclass{tufte-handout}

%\geometry{showframe}% for debugging purposes -- displays the margins

\usepackage{amsmath}
\usepackage{amssymb}
\usepackage{amsthm}
\usepackage{thmtools}
\declaretheorem[numberwithin=section]
{theorem, definition}
\declaretheorem{lemma, proposition, corollary}[
style=plain,
numberwithin=theorem
]
% Set up the images/graphics package
\usepackage{graphicx}
\setkeys{Gin}{width=\linewidth,totalheight=\textheight,keepaspectratio}
\graphicspath{{graphics/}}

\title{MATH 356 notes}
\author{Shan Gao}
\date{Fall 2023}

% The following package makes prettier tables.  We're all about the bling!
\usepackage{booktabs}

% The units package provides nice, non-stacked fractions and better spacing
% for units.
\usepackage{units}

% The fancyvrb package lets us customize the formatting of verbatim
% environments.  We use a slightly smaller font.
\usepackage{fancyvrb}
\fvset{fontsize=\normalsize}

% Small sections of multiple columns
\usepackage{multicol}

% Provides paragraphs of dummy text
\usepackage{lipsum}

% These commands are used to pretty-print LaTeX commands
\newcommand{\doccmd}[1]{\texttt{\textbackslash#1}}% command name -- adds backslash automatically
\newcommand{\docopt}[1]{\ensuremath{\langle}\textrm{\textit{#1}}\ensuremath{\rangle}}% optional command argument
\newcommand{\docarg}[1]{\textrm{\textit{#1}}}% (required) command argument
\newenvironment{docspec}{\begin{quote}\noindent}{\end{quote}}% command specification environment
\newcommand{\docenv}[1]{\textsf{#1}}% environment name
\newcommand{\docpkg}[1]{\texttt{#1}}% package name
\newcommand{\doccls}[1]{\texttt{#1}}% document class name
\newcommand{\docclsopt}[1]{\texttt{#1}}% document class option name

\begin{document}
	\maketitle
	\begin{definition}[Probability Space]
		A probability space is a triple $(\Omega, \mathcal{F}, \mathbb{P})$, where $\Omega$ is the sample space, $\mathcal{F}$ are the events subsets of $\Omega$ and $\mathbb{P}$ is the probability measure function, so $\mathbb{P} \colon \mathcal{F} \to \mathbb{R}$.  

		
	\end{definition}
	\begin{definition}[Kolmogorov Axioms]
		\begin{itemize}
			\item $ 0 \leq \mathbb{P}(A) \leq 1 \quad \forall A \in \mathcal{F}$ 
			\item $\mathbb{P}(\emptyset) = 0$
			\item If $\{A_i\}_{i \geq 1}$ are disjoint events, that means if $i \neq j$, then $A_i \cap A_j = \emptyset$, then 
				\[\mathbb{P} \left( \bigcup_{i \geq 1} A_i \right) = \sum_{i \geq1} \mathbb{P}(A_i)\]
		\end{itemize}	
	\end{definition}
	A common example is the roll of the die where $\Omega = \{1,2,3,4,5,6\}$, where $\mathbb{P}(i) = \frac{1}{6}$ where $i \in [6]$.
	We can model the roll of two dice, where $\Omega = [6] \times [6]$, then that means $\mathbb{P}\{(i,j)\} = \frac{1}{36}$ where $i,j \in \Omega$.
	Our probability measure function is not always uniform. For example if you decided to roll a die, that had some bias. Let's say it lands with probability $p$ on head. We know that $\Omega = \{0,1\}$. If we were to toss a coin three times, then $\Omega = \{0,1\}^3 = \{(i,j,k) \colon i,j,k \in \{0,1\}\}$, if we were to apply our measure function then 
	\[\mathbb{P} \left( \{i,j,k\} \right) = p^\text{\# coord $=1$ } \cdot (1-p)^\text{ \# coords $=0$}\]
	For example 
	\[\mathbb{P} \left( \{1,0,1\} \right) = p^2\cdot (1-p)\]
	Another example is if we tossed a coin infinitely many times, what would we get? 
	The sample space is $$\Omega = \{0,1\}^\mathbb{N} = \big\{\{b_i\}_{i \geq 1} \colon b_i \in \{0,1\} \text{ for each $i$ } \big\}$$
	To better visualise this, the sample space is uncountable, this is similar to if we picked any real number in the interval $[0,1]$. If we let $[a,b] \subseteq [0,1]$, where $0\leq a\leq b \leq 1$, then 
	\[\mathbb{P}([a,b]) = b-a \implies \forall x \in [0,1] \quad \mathbb{P}(x) = 0\]
	\begin{proof}
		Fix $\varepsilon > 0$, then $\mathbb{P}([x-\varepsilon, x + \varepsilon]) = 2\varepsilon$, but that means, $\mathbb{P}(\{x\}) < 2\varepsilon$, so we are done
	\end{proof}
	\section{lecture 2}
	\begin{corollary}[finite additivity]
	We can recall Kolmogorov's axioms, these axioms actually mean that for any finite sequence of events in $\mathcal{F}$ which are disjoint then 
	$$\mathbb{P}(\bigcup_{i = 1}^n A_i) = \sum_{i = 1}^n \mathbb{P}(A_i)$$
	The proof of this is trivial
	\end{corollary}
	A consequence of this is that $\mathbb{P}(A \cup A^c) = \mathbb{P}(A) + \mathbb{P}(A^c)$ for any event, so $\mathbb{P}(A) = 1- \mathbb{P}(A^c)$
	We now introduce some experiment models that could be helpful
	\begin{enumerate}
		\item If we are sampling with replacement, and the probability model, with a uniform probability measure, is 
			$$ \Omega = [50]\times [50] \times [50]$$
			then the probability of any point in the outcome is 
			$$\mathbb{P}(\text{outcome}) = \frac{1}{50^3}$$
		\item If we are sampling without replacement, then 
			$$ \Omega = \{(i,j,k) \colon i,j,k \in [50] \text{ and } i\neq j \neq k \} $$
			For example 
			$$ \mathbb{P}(\{22,7,48\}) = \frac{1}{50 \cdot 49 \cdot 48}$$

	\end{enumerate}

	
\end{document}
